\scanpage{20}
% archive sheet number: 19
% slanted
\margin{expressions en $(\varepsilon_0,\rho'_0)$}
\[k)\ \left\{\begin{aligned}
\varepsilon_1&=\rho'_0(\varepsilon_0)=\rho'_0\varepsilon_0{\rho'_0}^{-1}\\
\sigma&=\rho'_0\varepsilon_0^{-1}\\
\rho_0&=\varepsilon_0{\rho'_0}^{-1}\varepsilon_0\\
\rho_1&=\varepsilon_0^2{\rho'_0}^{-1}\\
\rho'_1&=\varepsilon_0^{\deleted{\unsure{-1}}}(\rho'_0)=\varepsilon_0\rho'_0\varepsilon_0^{-1}
\end{aligned}\right.\]
% slanted
\margin{expressions en $(\varepsilon_0,\rho_1)$}
\[\text{\textcircled{l})}\ \left\{\begin{aligned}
\varepsilon_1&=\rho_1^{-1}(\varepsilon_0)=\rho_1^{-1}\varepsilon_0\rho_1\\
\sigma&=\rho_1^{-1}\varepsilon_0\\
\rho_0&=\varepsilon_0^{-1}(\rho_1)=\varepsilon_0^{-1}\rho_1\varepsilon_0\\
\rho'_1&=\rho_1^{-2}\\
\rho'_0&=\varepsilon_0^{-1}(\rho_1^{-2})=\varepsilon_0^{-1}\rho_1^{\unsure{-1|-7|-2}}\varepsilon_0
\end{aligned}\right.\]
% slanted
\margin{expressions en $(\varepsilon_0,\rho'_1)$}
\[\text{\textcircled{\unsure{m}}}\ \left\{\begin{aligned}
\deleted{\unsure{\sigma=\varepsilon_0^{-1}\rho'_1}}\quad \varepsilon_1&={\rho'_1}^{-1}(\varepsilon_0)={\rho'_1}^{-1}\varepsilon_0\rho'_1\\
\sigma&=\varepsilon_0^{-1}\rho'_1\\
\rho_1&=\varepsilon_0{\rho'_1}^{-1}\varepsilon_0\\
\rho'_0&=\varepsilon_0^{-1}(\rho'_1)=\varepsilon_0^{-1}\rho'_1\varepsilon_0\\
\rho_0&={\rho'_1}^{-1}\varepsilon_0^2
\end{aligned}\right.\]
% slanted
\margin{19.6.83}

\uline{NB} parmi les sept éléments $\rho_i,\rho'_i,\varepsilon_i,\sigma$, \deleted{deux \unsure{q}} \inserted{une}
paire quelconque qui n'est pas formée uniquement
d'éléments parmi les 4 $\rho$ [ce qui fait $\binom{7}{2}-\binom{4}{2}$
$=21-6=15$ paires) engendrent $H$. D'autre part,
les quatre éléments $\unsure{\rho_i|\rho}$ $(\rho_0,\rho_1,\rho'_0,\rho'_1)$ n'engendrent
pas $H=\langle \varepsilon_0,\varepsilon_1\mid \varepsilon_0\varepsilon_1\varepsilon_0=\varepsilon_1\varepsilon_0\varepsilon_1\rangle$, \unsure{par} $\unsure{=|\simeq}$ mod commutateurs\unsure{...}.

Il faudrait aussi regarder l'extension canonique
de $\mathbb{Z}/2\mathbb{Z}$ par
\[\begin{aligned}
\unsure{S\mathcal{T}|S\mathscr{T}}_{1,1}^+&=\langle \varepsilon_0,\varepsilon_1\mid \varepsilon_0\varepsilon_1\varepsilon_0=\varepsilon_1\varepsilon_0\varepsilon_1\rangle\\
&\simeq\langle\rho,\sigma\mid \rho^{-3}=\sigma^2\rangle\\
&\simeq\langle\rho,\varepsilon_0\mid \rho_{\unsure{0|\vphantom{0}}}\varepsilon_0\rho_{\unsure{0|\vphantom{0}}}^{-1}\varepsilon_0\rho=1\rangle
\end{aligned}\]
donnant naissance (par \unsure{division} par la relation $\omega=1$)
à $\mathrm{Gl}(2,\mathbb{Z})$ ext. de $\mathbb{Z}/2\mathbb{Z}$ par $\mathrm{Sl}(2,\mathbb{Z})$.
